% Nouveau EGA V (cote 25): archive cote 25, site page 75 (PDF page 76).
% typescript. Draft compared: none.
% AI-transcribed by Claude Opus 5.5 (claude-opus-5-5, adaptive thinking, effort high), Message Batches API, batch ega25b (msgbatch_01A7n5YuXubUPNK1zsepi8RG), 2026-11-13; checked against both supplied scan parts; see CHECK_DECISIONS.tsv.
\scanpage{25/75}
\mspage{70}

20.13. \uline{Morphismes élémentaires et théorème de} M. \uline{Artin}.

\uline{Définition} 13.1. Un morphisme $f:X\hand{\to} Y$ de préschémas est dit ``morphisme
\var{\xout{élém}}{}{typewriter strike-out}{ega25_25_075-v0001} élémentaire'' si X est \var{\xout{un}}{}{typewriter strike-out}{ega25_25_075-v0002} Y-isomorphe à un préschéma de la forme
$X' - Z$ , où $X'\hand{\to} X$ est un morphisme lisse, projectif, à fibres géomé-%
triques connexes\var{\xout{\unsure{,}}}{}{typewriter strike-out}{ega25_25_075-v0003} de dimension 1, et où Z est un sous-préschéma fermé
de X', \var{\xout{étale sur}}{}{typewriter strike-out}{ega25_25_075-v0004} tel que le morphisme $Z\hand{\to} Y$ soit étale, surjectif, et
de degré constant. Un morphisme est dit \var{\xout{\unsure{multip-élémentaire}}}{}{typewriter strike-out}{ega25_25_075-v0005} poly-élémen-%
taire s'il est un composé de morphismes élémentaires. Un préschéma X
sur un corps k est dit polyélementaire (sur k), si le morphisme structu-%
ral $X\hand{\to} \mathrm{Spec}(k)$ est poly-élémentaire.

\uline{Théorème} 13.2. \var{\inserted[typed]{(M.Artin)}}{(M.Artin)}{insertion}{ega25_25_075-v0006} Soient X un préschéma \var{\inserted[typed]{géom. irréd.}}{géom. irréd.}{insertion}{ega25_25_075-v0007} sur un corps k \var{\inserted[typed]{parfait infini}}{parfait infini}{insertion}{ega25_25_075-v0008}, x un point lisse
de X, alors x admet un système fondamental de voisinages ouverts
polyélémentaires.

Quitte à remplacer X par un voisinage donné de x, \var{\deleted{on \var{\xout{\unsure{peut suppo}}}{}{typewriter strike-out}{ega25_25_075-v0010} es}}{}{deletion}{ega25_25_075-v0009} \var{\inserted{il revient au}}{il revient au}{insertion}{ega25_25_075-v0011}
\var{\deleted{ramené à}}{}{deletion}{ega25_25_075-v0012} \var{\inserted{même de}}{même de}{insertion}{ega25_25_075-v0013} montrer qu'il existe un voisinage ouvert élémentaire de x dans X.
\var{\deleted{On peut supposer évidemment X lisse, et}}{}{deletion}{ega25_25_075-v0014} \var{\inserted[typed]{\var{\deleted{et affine}}{}{deletion}{ega25_25_075-v0016}}}{\var{\deleted{et affine}}{}{deletion}{ega25_25_075-v0016}}{insertion}{ega25_25_075-v0015} \var{\over{r}{R}}{R}{writing over}{ega25_25_075-v0017}aisonnant par récurrence sur
la dimension n de X, on \var{\xout{peut supposer}}{}{typewriter strike-out}{ega25_25_075-v0018} est ramené à montrer \var{\inserted[typed]{que si $n\hand{>}0$ ,}}{que si $n\hand{>}0$ ,}{insertion}{ega25_25_075-v0019} \var{\deleted{qu'}}{}{deletion}{ega25_25_075-v0020}il existe
un voisinage ouvert U de x et un morphisme élémentaire $f: U\hand{\to} V$,
V étant un schéma lisse sur k (nécessairement géom. irréd \var{\inserted[typed]{et}}{et}{insertion}{ega25_25_075-v0021} de dim n-1).
\var{\xout{Lorsque n=1}}{}{typewriter strike-out}{ega25_25_075-v0022} (Le cas n=0 est évidemment trivial, étant donné qu\var{\over{e}{'}}{'}{writing over}{ega25_25_075-v0023} alors
X est k-isomorphe à Spec(k), qui est polyélémentaire sur Spec(k), étant
entendu dans 13.1. qu'on n'exc\hand{l}ut pas le composé de la famille vide de
morphismes - on aurait dû le signaler sous une forme ou une autre dans
13.1. ). La nécessité de supposer k parfait au préalable apparait déjà
dans le cas n=1: alors on prend pour X' le modèle projectif normal
canonique du corps des fonctions K de X (cf Chap II) par.7), le fait que

\mspage{75}
